\subsection{Commercial BCA Database Creation} \label{subsec:commercial-bca-database}

After source-specific parsing, the records are normalised and checked against curated microbe and pathogen reference data.
The normalisation step removes formatting differences while retaining meaningful qualifiers.
Missing values stay missing, and failed reference matches receive an explicit review reason.
This keeps unresolved cases visible instead of turning them into assumed biological relationships.

The database creation stage classifies records into three evidence states.
A record enters \textit{final} when it contains BCA evidence and has no unresolved review reason.
A record enters \textit{review} when it contains BCA evidence but has an unresolved reference or crop--pathogen pair.
All remaining records enter \textit{discarded}.
These are authorisation evidence classes, not biological efficacy labels: a discarded record means that the pipeline found no BCA signal, not that the product or microorganism is ineffective.

The resulting commercial BCA database preserves the crop, pathogen, product, microbial active substance, source, authorisation, product function, and review fields available in each record.
Product names remain linked to the standardized microbial active substance value as traceable evidence.
Only rows with complete scientific crop and pathogen pairs enter the joined modelling table.
This prevents a product catalogue from being interpreted as a verified set of crop--pathogen uses.

The current ranking database is built from accepted Portuguese records because they provide the complete fields needed for crop--pathogen evaluation.
The other country outputs remain available for catalogue inspection and audit until their missing links can be validated.
The final and review classes also provide weak labels for the diagnostic triage model evaluated later in the paper.
